\cotefile{156-9}
\scanpage{50}
\mspage{46}
% archive sheet number: 49
On trouve donc que la structure d'\uline{atelier spatial} peut se décrire par les données
\[
(\unsure{\mathcal A|\mathscr A},\trianglelefteq,\mathrel{|\unsure{o|\circ}|})
\]
avec les seuls axiomes

\uline{At spat 1}\quad a) $\trianglelefteq$ est une relation d'ordre.

\hspace*{4em}b) $\mathrel{|o|}$ est symétrique \unsure{antiréflexive}

\uline{At spat 2}\quad Si $X\triangleleft Y$, on a $X\mathrel{|o|}Y$.

\uline{At spat 3}\quad L'application
\[
\begin{aligned}
\mathcal A &\longrightarrow \deleted{\unsure{\mathfrak S|S}}\,\mathrm{Fig}\,\text{él}\,(\unsure{P|\mathcal P}(\mathcal A))\\
X &\longmapsto \text{Dép}\,X
 = \bigl(\{X'^{\unsure{\circ|o}}\mid X'\unsure{\mathrel{\ooalign{\hfil$\trianglelefteq$\hfil\cr\hfil\raisebox{0.32ex}{\scriptsize$-$}\hfil\cr}}|\trianglelefteq}X\},\,\deleted{\text{\unsure{avec}}\,\unsure{\leq}}\,+\\
 &\hspace{2.6em}\deleted{\text{où }X'^{\circ}=\{Y}\quad
 \inserted{\begin{aligned}
 &\text{relation d'ordre sur Dép }\unsure{X|X'}\\[-2pt]
 &\text{déduite de celle de}\\[-2pt]
 &\widetilde X=\{x'\in\mathcal A\mid x\trianglelefteq X\}
 \end{aligned}}\bigr)
\end{aligned}
\]
est \uline{injective}. (\unsure{On|on} a \unsure{a} posé, pour $Y\in\mathcal A$,
\[
Y^{\circ}=\{Z\in\mathcal A\mid\forall Z'\in\mathcal A,\ Z'\mathrel{|o|}Y\Rightarrow Z'\mathrel{|o|}Z\}
\]

Je vais essayer de \unsure{reformuler} la structure, avec $\Sigma=\Sigma_{\mathcal A}$ comme ensemble de base.
\margin{$(\Sigma,$}

\uline{Définition} On appelle \uline{\unsure{spatios|spatio}} (ensembliste) un \deleted{\unsure{ensemble ordonné}}\inserted{triple} $(\Sigma,\leq,\deleted{\unsure{f|\ell|\rho}}C)$, où $\Sigma$ est un ensemble, $\leq$ une relation d'ordre sur $\Sigma$, $C$ une application de $\Sigma$ en lui-même
